% Part: normal-modal-logic
% Chapter: filtrations
% Section: finite

\documentclass[../../../include/open-logic-section]{subfiles}

\begin{document}

\olfileid{nml}{fil}{fin}

\olsection{ফিল্ট্রেশন সসীম}

উপ-!!{formula}s-এর অধীনে বদ্ধ যেকোনো সমষ্টি~$\Gamma$-র জন্য আমরা ফিল্ট্রেশন সংজ্ঞায়িত করেছি। সংজ্ঞা একাই ফিল্ট্রেশন সসীম হওয়া নিশ্চিত করে না। বাস্তবে $\Gamma$ অসীম হলে (যেমন সব !!{formula}s-এর সমষ্টি হলে) ফিল্ট্রেশনও অসীম হতে পারে। কিন্তু $\Gamma$ সসীম হলে (যেমন প্রদত্ত !!{formula}~$!A$-র উপ-!!{formula}s-এর সমষ্টি হলে), $\Gamma$-র মধ্য দিয়ে যেকোনো ফিল্ট্রেশনও সসীম।

\begin{prop}\ollabel{prop:filt-are-finite}
  $\Gamma$ সসীম হলে, $\Gamma$-র মধ্য দিয়ে মডেল~$\mModel{M}$-এর যেকোনো ফিল্ট্রেশন $\mModel{M^*}$-ও সসীম।
\end{prop}

\begin{proof}
  $W^*$-র আকার হলো সমতুল্যতা-সম্বন্ধ~$\equiv$-এর অধীনে ভিন্ন ভিন্ন শ্রেণি~$[w]$-এর সংখ্যা। একই শ্রেণিতে থাকা যেকোনো দুটি জগৎ $u$, $v$—অর্থাৎ, যেসব $u$ ও $v$-র জন্য~$u \equiv v$—$\Gamma$-র সব !!{formula}s~$!A$-তে একমত: $!A \in \Gamma$ হলে হয় $u$ ও $v$ উভয়েই $!A$ সত্য করে, নয়তো কেউই করে না। তাই প্রতিটি শ্রেণি~$[w]$ $\Gamma$-র একটি উপসমষ্টির সঙ্গে সম্পর্কিত; সেটি হলো~$[w]$-র জগৎগুলিতে সত্য সব $!A \in
  \Gamma$-র সমষ্টি। ভিন্ন দুটি শ্রেণি $[u]$ ও $[v]$ $\Gamma$-র একই উপসমষ্টির সঙ্গে সম্পর্কিত হতে পারে না। কারণ $u$ ও $v$-তে সত্য !!{formula}s-এর সমষ্টি একই হলে, এই $u$ ও $v$ $\Gamma$-র সব !!{formula}s-এ একমত; অর্থাৎ $u \equiv v$। কিন্তু তাহলে $[u] = [v]$। সুতরাং $W^*$ থেকে $\Pow{\Gamma}$-তে একটি !!a{injective} অপেক্ষক আছে, ফলে $\card{W^*} \le \card{\Pow{\Gamma}}$। তাই সমষ্টিটির প্রতিটি $!A$-র সত্যমান বিবেচনায়, $\Gamma$-তে $n$টি বাক্য থাকলে $W^*$-র কার্ডিনালিটি~$2^n$-এর বেশি নয়।
\end{proof}

\end{document}
